\section{Related Work}
\label{sec:related}

Three lines of work meet in this study: attempts to make networks interpretable
by changing their architecture, tensor-factorised Mixture-of-Experts layers,
and the mechanistic analysis of grokking on modular arithmetic. They have
developed largely independently, and part of the contribution here is to bring
them into contact.

\subsection{Polysemanticity and architectural remedies}

The difficulty that motivates all of this is that individual neurons do not
correspond to individual features. \citet{elhage2022superposition} give the
standard account: when a network must represent more features than it has
dimensions, it stores them in superposition, and each neuron responds to
several unrelated inputs. Interpretability then requires disentangling
features from neurons rather than reading neurons directly.

Two families of response exist. The first leaves the model alone and analyses
it afterwards, most prominently by training a sparse autoencoder on the
activations to recover an overcomplete dictionary of features
\citep{bricken2023monosemanticity}. This works on any pretrained model, but it
introduces a second learned object whose own faithfulness must then be
established.

The second family changes the architecture so that the units are monosemantic
by construction. \citet{park2025monet} are the fullest development of this
idea. Their MONET layer replaces the feed-forward block with a very large
mixture of experts and, crucially, does not store the experts separately: a
layer is partitioned into $\sqrt{N}$ \emph{bottom} factors and $\sqrt{N}$
\emph{top} factors, and expert $(i,j)$ is composed on demand from bottom
factor $i$ and top factor $j$. Parameter count then scales as $\sqrt{N}$
rather than $N$, which lets them instantiate $262\,144$ experts per layer.
They report that experts specialise by domain, and demonstrate control by
ablating experts associated with a programming language or with toxic content.

Their evidence is correlational by construction. An expert is assigned to a
domain when its routing score on that domain exceeds its score elsewhere by a
factor of two, and toxicity experts are located by correlating routing scores
with a toxicity label. The authors state plainly in their limitations that
these selection criteria are rudimentary and that quantitative evaluation of
automated interpretability remains open. The obstacle is not one of effort: on
natural language there is no independent account of which features the model
ought to represent, so a specialisation claim can be illustrated but not
verified. \citet{bereska2024mechinterp}, surveying the field, make the same
point from the other direction and call explicitly for algorithmic testbeds
with known ground truth against which interpretability methods can be scored.

\subsection{Factorised and large-scale mixtures of experts}

Mixture-of-Experts layers were introduced to decouple parameter count from
per-token compute \citep{shazeer2017moe}, with interpretability arriving much
later as a secondary motivation. Scaling the number of experts runs into
routing cost, since a naive softmax over $N$ experts is linear in $N$.
\citet{lample2019productkeys} solve this with product keys: the query is split
in half, each half is matched against $\sqrt{N}$ sub-keys, and the top-$k$
over the Cartesian product is recovered from the top-$k$ of each half, giving
$O(\sqrt{N})$ routing. \citet{he2024peer} combines product-key routing with
single-neuron experts to reach a million experts per layer. In both, the
experts themselves remain independent parameters; only the routing is
factorised.

\citet{oldfield2024mumoe} take the further step of factorising the experts.
They observe that the parameters of $N$ linear experts form a third-order
tensor in $\R^{d_{\mathrm{out}} \times d_{\mathrm{in}} \times N}$, and replace
it with a CP or Tucker factorisation, so that the layer computes a multilinear
function of the input and the expert-coefficient vector without ever
materialising the individual experts. They report that expert specialisation
increases with expert count, measuring it by the drop in per-class accuracy
when an expert is ablated. MONET criticises this line for parameter counts
that still grow linearly in $N$ at fixed rank, and positions product-key
composition as the more economical alternative.

This disagreement is the starting point of the present work. Product-key
composition and a CP or Tucker factorisation are both factorisations of the
same expert tensor; they differ in how much structure is shared across
experts, and the literature compares them on cost rather than on what they do
to the function the network learns.

\subsection{Grokking and modular arithmetic}

\citet{power2022grokking} observed that small networks trained on algorithmic
tasks fit the training set almost immediately while test accuracy remains at
chance for thousands of further steps, then rises sharply. The delay makes the
transition from memorisation to generalisation externally visible as a single
number, which is what we exploit.

For modular addition the generalising solution has been reverse-engineered in
full. \citet{nanda2023progress} show that a one-layer transformer trained on
$c = (a+b) \bmod P$ represents the operands through a small set of Fourier
frequencies, multiplies the corresponding trigonometric terms through the
attention and MLP blocks, and reads out the logit for $c$ as a sum of
$\cos(\omega(a+b-c))$ terms; they use progress measures derived from this
structure to show that the algorithm forms gradually, well before test
accuracy moves. \citet{bussmann2023modadd} gives the corresponding analysis
for a plain MLP, where each hidden neuron approximates
$\mathrm{ReLU}\!\left(\cos(\omega a + \varphi_1) + \cos(\omega b +
\varphi_2)\right)$ and its outgoing weights follow $\cos(\omega c + \varphi_1
+ \varphi_2)$; notably, the MLP spreads its representation over a broad range
of frequencies rather than selecting a few, in contrast to the transformer.
\citet{chughtai2023universality} extend the picture to general group
composition and find only weak universality across seeds, which is the main
reason we report multiple seeds rather than a single run.

What this literature provides, and what is missing everywhere else in the
interpretability of expert layers, is an answer key. The features a
generalising network must represent are known exactly and there are only
$(P-1)/2$ of them, so the question of whether a unit is monosemantic has a
definite answer.

\subsection{Positioning}

Against this background our study occupies an otherwise empty position. The
MoE-interpretability literature evaluates specialisation on natural data
without ground truth; the tensor-decomposition literature compares
factorisations on parameter count and downstream accuracy; the grokking
literature analyses dense architectures. We hold the task, the data and the
optimisation protocol fixed, vary only the factorisation of the expert tensor,
and measure both when the generalising algorithm is found and what its learned
representation looks like in the basis the task itself defines. To our
knowledge the relationship between expert-tensor factorisation and the
memorisation-to-generalisation transition has not previously been measured.
